% GENERATED FILE. Edit canonical lessons, then run npm run generate:books.
% canonical-source-hash: fnv1a64:587b02c53fb07fc9
% canonical-lessons: SA-C121-vadyam, SA-C121-vina, SA-C121-pattah, SA-C121-ratnam, SA-C121-suvarnam

\chapter{Musical instrument, Veena, Cloth, Jewel, Gold}
\label{ch:sa-musical-instrument-veena-cloth-jewel-gold}
\hlchaptermodality{121}

\begin{chapteropening}
\textbf{By the end of this chapter:} \emph{I can say a musical instrument, a veena, a cloth, a jewel and gold.}

Complete the last lesson of chapter 121: i can say a musical instrument, a veena, a cloth, a jewel and gold.
\end{chapteropening}

\section[\texorpdfstring{\sk{वाद्यम्} (vādyam)}{वाद्यम् (vādyam)}]{\sk{वाद्यम्} (vādyam) — a musical instrument}

\label{lesson:SA-C121-vadyam}



\begin{quote}
Before the new one: say the Sanskrit for a book, then the Sanskrit for a picture.
\end{quote}

\subsection*{You'll want to know: \sk{वाद्यम्}}
\textbf{\sk{वाद्यम्}} — \emph{vādyam} — \textquotedblleft{}a musical instrument\textquotedblright{}.

\begin{grammarlens}[title={What you've built}]
One more word for this chapter.
\end{grammarlens}

\subsection*{Your turn}
\begin{itemize}
\raggedright
  \item \emph{Say it:} \emph{vādyam}
  \item \emph{Say it:} \emph{vādyam}, once more
  \item \emph{Say it:} say \emph{vādyam}
  \item \emph{Recall:} say the Sanskrit for a book, then the Sanskrit for a picture, then say \emph{vādyam} again
\end{itemize}

\begin{tcolorbox}[breakable,colback=teal!4,colframe=teal!35!black,title={Before you move on}]
What does \sk{वाद्यम्} mean? (\textquotedblleft{}A musical instrument\textquotedblright{}.) Say it once more.
\end{tcolorbox}
\section[\texorpdfstring{\sk{वीणा} (vīṇā)}{वीणा (vīṇā)}]{\sk{वीणा} (vīṇā) — a veena}

\label{lesson:SA-C121-vina}



\begin{quote}
Before the new one: say the Sanskrit for a picture, then the Sanskrit for a musical instrument.
\end{quote}

\subsection*{You'll want to know: \sk{वीणा}}
\textbf{\sk{वीणा}} — \emph{vīṇā} — \textquotedblleft{}a veena\textquotedblright{}.

\begin{grammarlens}[title={What you've built}]
One more word for this chapter.
\end{grammarlens}

\subsection*{Your turn}
\begin{itemize}
\raggedright
  \item \emph{Say it:} \emph{vīṇā}
  \item \emph{Say it:} \emph{vīṇā}, once more
  \item \emph{Say it:} say \emph{vīṇā}
  \item \emph{Recall:} say the Sanskrit for a picture, then the Sanskrit for a musical instrument, then say \emph{vīṇā} again
\end{itemize}

\begin{tcolorbox}[breakable,colback=teal!4,colframe=teal!35!black,title={Before you move on}]
What does \sk{वीणा} mean? (\textquotedblleft{}A veena\textquotedblright{}.) Say it once more.
\end{tcolorbox}
\section[\texorpdfstring{\sk{पट्टः} (paṭṭaḥ)}{पट्टः (paṭṭaḥ)}]{\sk{पट्टः} (paṭṭaḥ) — a cloth, a strip}

\label{lesson:SA-C121-pattah}



\begin{quote}
Before the new one: say the Sanskrit for a musical instrument, then the Sanskrit for a veena.
\end{quote}

\subsection*{You'll want to know: \sk{पट्टः}}
\textbf{\sk{पट्टः}} — \emph{paṭṭaḥ} — \textquotedblleft{}a cloth, a strip\textquotedblright{}.

\begin{grammarlens}[title={What you've built}]
One more word for this chapter.
\end{grammarlens}

\subsection*{Your turn}
\begin{itemize}
\raggedright
  \item \emph{Say it:} \emph{paṭṭaḥ}
  \item \emph{Say it:} \emph{paṭṭaḥ}, once more
  \item \emph{Say it:} say \emph{paṭṭaḥ}
  \item \emph{Recall:} say the Sanskrit for a musical instrument, then the Sanskrit for a veena, then say \emph{paṭṭaḥ} again
\end{itemize}

\begin{tcolorbox}[breakable,colback=teal!4,colframe=teal!35!black,title={Before you move on}]
What does \sk{पट्टः} mean? (\textquotedblleft{}A cloth, a strip\textquotedblright{}.) Say it once more.
\end{tcolorbox}
\section[\texorpdfstring{\sk{रत्नम्} (ratnam)}{रत्नम् (ratnam)}]{\sk{रत्नम्} (ratnam) — a jewel}

\label{lesson:SA-C121-ratnam}



\begin{quote}
Before the new one: say the Sanskrit for a veena, then the Sanskrit for a cloth.
\end{quote}

\subsection*{You'll want to know: \sk{रत्नम्}}
\textbf{\sk{रत्नम्}} — \emph{ratnam} — \textquotedblleft{}a jewel\textquotedblright{}.

\begin{grammarlens}[title={What you've built}]
One more word for this chapter.
\end{grammarlens}

\subsection*{Your turn}
\begin{itemize}
\raggedright
  \item \emph{Say it:} \emph{ratnam}
  \item \emph{Say it:} \emph{ratnam}, once more
  \item \emph{Say it:} say \emph{ratnam}
  \item \emph{Recall:} say the Sanskrit for a veena, then the Sanskrit for a cloth, then say \emph{ratnam} again
\end{itemize}

\begin{tcolorbox}[breakable,colback=teal!4,colframe=teal!35!black,title={Before you move on}]
What does \sk{रत्नम्} mean? (\textquotedblleft{}A jewel\textquotedblright{}.) Say it once more.
\end{tcolorbox}
\section[\texorpdfstring{\sk{सुवर्णम्} (suvarṇam)}{सुवर्णम् (suvarṇam)}]{\sk{सुवर्णम्} (suvarṇam) — gold}

\label{lesson:SA-C121-suvarnam}



\begin{quote}
Before the new one: say the Sanskrit for a cloth, then the Sanskrit for a jewel.
\end{quote}

\subsection*{You'll want to know: \sk{सुवर्णम्}}
\textbf{\sk{सुवर्णम्}} — \emph{suvarṇam} — \textquotedblleft{}gold\textquotedblright{}.

\begin{grammarlens}[title={What you've built}]
One more word for this chapter.
\end{grammarlens}

\subsection*{Your turn}
\begin{itemize}
\raggedright
  \item \emph{Say it:} \emph{suvarṇam}
  \item \emph{Say it:} \emph{suvarṇam}, once more
  \item \emph{Say it:} say \emph{suvarṇam}
  \item \emph{Recall:} say the Sanskrit for a cloth, then the Sanskrit for a jewel, then say \emph{suvarṇam} again
\end{itemize}

\begin{tcolorbox}[breakable,colback=teal!4,colframe=teal!35!black,title={Before you move on}]
What does \sk{सुवर्णम्} mean? (\textquotedblleft{}Gold\textquotedblright{}.) Say it once more.
\end{tcolorbox}
